% Grandissement et théorème de Thalès : le rayon passant par O (exercice 11, chapitre 18)
\begin{tikzpicture}[x=1.4cm,y=1.4cm]
\draw[otInk!60,line width=0.5pt] (-4.4,0) -- (3.2,0);
\draw[otNavy,line width=1.2pt,{Stealth[length=2.2mm]}-{Stealth[length=2.2mm]}] (0,-1.1) -- (0,1.4);
\solution{
  \fill[otNavy!12] (-4,0) -- (-4,1.2) -- (0,0) -- cycle;
  \fill[otGreen!18] (2.4,0) -- (2.4,-0.72) -- (0,0) -- cycle;}
\draw[otNavy,line width=1.2pt,-{Stealth[length=2.2mm]}] (-4,0) -- (-4,1.2) node[pcLegende,anchor=south,text=otNavy] {B};
\node[pcLegende,anchor=north] at (-4,-0.03) {A};
\draw[otGreenDark,line width=1.2pt,-{Stealth[length=2.2mm]}] (2.4,0) -- (2.4,-0.72) node[pcLegende,anchor=north,text=otGreenDark] {B'};
\node[pcLegende,anchor=south] at (2.4,0.03) {A'};
\node[pcLegende,anchor=north east] at (0,-0.02) {O};
\draw[otOrange!85!black,line width=0.8pt] (-4,1.2) -- (2.4,-0.72);
\solution{\node[pcLegende,anchor=west,align=left] at (2.6,0.5) {triangles OAB et OA'B' :\\les droites (AB) et (A'B')\\sont parallèles (Thalès)};}
\end{tikzpicture}
\par\smallskip\legende{Le rayon issu de B qui passe par le centre optique O n'est pas dévié et passe par B'.}
